\section{Problem and Semantic Boundaries}
\label{sec:problem}
\subsection{Selective quantum opportunity analysis}
Let $P$ be a classical program, $D$ its permitted inputs, and $C$ its observable
contract. We write $\operatorname{Obs}(P,x)$ for the returned value, relevant
effects, exceptions, and ordering observable on input $x$. A candidate region
$r$ is a source region together with the definitions and uses required to
interpret it. A mapping proposal $M$ describes an encoding, a search predicate
or optimization objective, a quantum procedure, decoding, and integration
conditions. A future transformed program $P[M/r]$ would need to satisfy
\begin{equation}
  \forall x\in D:\quad
  \operatorname{Obs}(P[M/r],x)\;\mathcal{R}_{C}\;\operatorname{Obs}(P,x),
  \label{eq:contract}
\end{equation}
where $\mathcal{R}_{C}$ is the relation permitted by the original contract.
For deterministic exact outputs this is equality of the required observations.
A probabilistic or approximate relation is admissible only when the contract
explicitly permits it. No tolerance is inferred from the use of a quantum solver.
Equation~\ref{eq:contract} states the eventual correctness requirement; passing
the local checks below does not prove it.

The analysis output is a tuple $(r,e,f,M,a,U)$: candidate regions $r$, structural
eligibility $e$, formulation family $f$, conditional mapping $M$, adoption
recommendation $a$, and unresolved obligations $U$. Eligibility takes values
$\{\yes,\no,\unknown\}$. A \yes{} requires a concrete correspondence between
the computation and a supported formulation, rather than only an algorithm name.
Adoption depends on additional implementation and resource evidence. Thus
$e=\yes$ and $a=\classical$ can be consistent. A conditional plan is still
required for a structural \yes{} in the prospective protocol.

\begin{table}[tb]
\caption{Evidence required at different stages. Later stages are not inferred
from earlier ones. Current experiments primarily concern the first four rows.}
\label{tab:stages}
\begin{tabularx}{\linewidth}{lYY}
\toprule
Stage & Required output & Distinct unresolved question \\
\midrule
WHERE & Candidate spans and dependencies & Is this an admissible region? \\
WHETHER & Mapping evidence and conditions & Is adoption worthwhile? \\
WHAT/HOW & Intent, encoding, predicate/objective & Does the formulation match the core? \\
PLAN & Decoding and obligation disposition & Is the full contract accounted for? \\
REFACTOR & Implemented hybrid replacement & Does integration preserve behavior? \\
VALIDATE & Semantic and resource evidence & Is correctness or benefit established? \\
\bottomrule
\end{tabularx}
\end{table}

The historical benchmark combines some of these conceptual stages in serialized
task identifiers. We retain those formats when replaying historical material and
use the named concepts in this paper to avoid conflating protocol versions.
Candidate nomination is also distinct from suitability: a validation scan may
be expressible as search yet offer no useful offloading opportunity. Whether it
belongs in the intended candidate population requires an explicit reference
policy, not an after-the-fact penalty for a model's choice.

\subsection{Supported formulations}
For search, the model must define a finite domain $\Omega(x)$ and a predicate
$g_x(z)$, then explain how marked states correspond to the source computation.
Grover's algorithm supplies a search primitive under oracle
assumptions~\cite{grover}. It does not itself supply the reversible oracle,
input-loading mechanism, mandated first witness, or no-solution certificate.
Those remain explicit plan obligations.

For optimization, a QUBO has the form
\begin{equation}
 E_x(z)=c_x+\sum_i h_i(x)z_i+\sum_{i<j}J_{ij}(x)z_i z_j,
 \qquad z\in\{0,1\}^{n_x}.
 \label{eq:qubo}
\end{equation}
The proposal must relate variables to source choices, encode constraints, specify
optimization direction, and decode auxiliary variables. Ising formulations
provide established constructions for combinatorial problems~\cite{lucas}.
QAOA is a possible approximate optimization procedure~\cite{qaoa}; naming it
does not certify exact global optimality or preserve a deterministic tie rule.
The current positive scope excludes general linear-system and estimation families.
Numerical programs can still contain search subregions, which are assessed locally.

\subsection{A running contract example}
Consider a vertex-cover routine that returns a minimum-cardinality cover and,
among ties, the lexicographically smallest tuple of selected vertex indices.
Its surrounding program turns the tuple into an ordered work list and a change
report. Preserving cardinality without preserving the tuple can therefore alter
the external report. Numeric bitmask order and tuple order are different
relations and must not be substituted silently.

Let $U(z)=\sum_{(i,j)\in E}(1-z_i)(1-z_j)$ count uncovered edges and let
$A=2^n$. One illustrative exact formulation is
\begin{equation}
 E(z)=\sum_{i=0}^{n-1}\bigl(A-2^{n-1-i}\bigr)z_i
       +P U(z),\qquad P=nA+1.
 \label{eq:cover}
\end{equation}
Every selection cost is positive and their total is less than $nA$. Consequently
an infeasible assignment is more expensive than the all-selected feasible
assignment. Among feasible assignments, the range of the binary-weight tie term
is less than $A$, so cardinality dominates. Among assignments of equal cardinality,
the first differing index dominates all later tie terms and favors inclusion of
the earlier vertex. This illustrates why semantic obligations determine the
objective. It is not an experimental result, a required reference formula, or a
claim that its coefficients are practical on a particular device. Equivalent
encodings and smaller valid penalties must remain admissible.

\figplaceholder{contract-example}{2.1cm}{Show the classical core, minimum-cardinality
selection, tuple ordering, and downstream work-list/report dependencies. Highlight
where a bitmask-based tie rule changes an observable output.}
{Illustrative contract dependencies for the vertex-cover case. The figure is
reserved for a source-linked explanation rather than experimental measurements.}
